%!TEX TS-program = xelatex
\documentclass[]{../../../../friggeri-cv}

% Variant source:
%   Anthropic - Engineering Manager, Connectivity
% This file is tailored from ../../../../cv.tex and intentionally does not
% replace the master CV source.

% Packages
\usepackage{pifont}
\usepackage{graphicx}
\usepackage{hyperref}
\usepackage{marvosym}
\usepackage{xcolor}
\usepackage{enumitem}
\usepackage{fancyhdr}

\setmainfont
[Mapping=tex-text, Color=textcolor, Ligatures=NoCommon,
BoldFont=texgyreheros-bold.otf,
ItalicFont=texgyreheros-italic.otf,
BoldItalicFont=texgyreheros-bolditalic.otf
]
{texgyreheros-regular.otf}
\newfontfamily\cvbodyfont[Ligatures=NoCommon]{Helvetica Neue}
\newfontfamily\cvthinfont[Ligatures=NoCommon]{Helvetica Neue UltraLight}
\newfontfamily\cvheadingfont[Ligatures=NoCommon]{Helvetica Neue Condensed Bold}
\renewcommand{\bodyfont}{\cvbodyfont}
\renewcommand{\thinfont}{\cvthinfont}
\renewcommand{\headingfont}{\cvheadingfont}

\graphicspath{ {../../../../images/} }
\definecolor{linkaccent}{HTML}{1E6291}
\definecolor{profileweb}{HTML}{B45309}
\definecolor{profilelinkedin}{HTML}{0A66C2}
\definecolor{profilegithub}{HTML}{24292F}
\definecolor{profilescholar}{HTML}{4285F4}
\newcommand{\profilelink}[3]{{\hypersetup{urlcolor=#1}\addfontfeature{Color=#1}\href{#2}{#3}}}
\hypersetup{
    colorlinks=true,
    urlcolor=linkaccent,
    linkcolor=linkaccent,
    citecolor=linkaccent,
    pdfborder={0 0 0},
    pdfauthor={Christos Hadjinikolis},
    pdftitle={Christos Hadjinikolis CV - Anthropic Engineering Manager Connectivity}
}

% Variant-scoped header override: restores a compact charcoal band while
% keeping a real space between first and last name for PDF text extraction.
\renewcommand{\header}[5]{%
  \begin{tikzpicture}[remember picture,overlay]
    \node [rectangle, fill=darkgray, anchor=north, minimum width=\paperwidth, minimum height=3.75cm] (box) at (current page.north){};
    \node [anchor=north] (name) at ([yshift=-0.72cm] box.north) {%
      \fontsize{36pt}{50pt}\color{white}%
      {\thinfont #1\ }{\bodyfont #2}%
    };
    \node [anchor=north] at ([yshift=-0.05cm] name.south) {%
      \fontsize{13pt}{20pt}\color{white}%
      \thinfont #3%
    };
    \node [anchor=north] at ([yshift=-0.55cm] name.south) {%
      \fontsize{11pt}{18pt}\color{white}%
      \hypersetup{urlcolor=white,linkcolor=white}\thinfont #4%
    };
    \node [anchor=north] at ([yshift=-3.08cm] box.north) {%
      \fontsize{10.5pt}{16pt}\color{white}%
      \hypersetup{urlcolor=white,linkcolor=white}\thinfont {\Letter\hspace{0.2em} #5}%
    };
  \end{tikzpicture}
  \vspace{2.2cm}
  \vspace{-2\parskip}
}

\renewcommand{\entry}[4]{%
  \parbox[t]{\dimexpr\textwidth-\entryleftinset\relax}{%
    \textbf{#2}%
    {\footnotesize\addfontfeature{Color=lightgray}\hspace{0.6em}#1}%
    \hfill%
    {\small\color{textcolor} #3}\\%
    #4\vspace{\parsep}%
  }\\}

\renewcommand{\entrytwo}[4]{%
  #1&\parbox[t]{14cm}{%
    \textbf{#2}%
    \hfill%
    {\small\color{textcolor} #3}\\%
    #4\vspace{\parsep}%
  }\\}

% Set up footer
\pagestyle{fancy}
\fancyhf{}  % Clear header/footer
\fancyfoot[C]{\small Page \thepage}
\renewcommand{\headrulewidth}{0pt}
\renewcommand{\footrulewidth}{0pt}

% XeLaTeX image fix
\ifxetex
  \usepackage{letltxmacro}
  \setlength{\XeTeXLinkMargin}{1pt}
  \LetLtxMacro\SavedIncludeGraphics\includegraphics
  \def\includegraphics#1#{%
    \IncludeGraphicsAux{#1}%
  }%
  \newcommand*{\IncludeGraphicsAux}[2]{%
    \XeTeXLinkBox{%
      \SavedIncludeGraphics#1{#2}%
    }%
  }%
\fi

\begin{document}
\header
    {Christos}
    {Hadjinikolis}
    {Engineering Manager | AI Connectivity \& Platform Reliability}
    {Ph.D. Computer Science | UCL Research Associate | AI Standards Expert @ \href{https://www.cencenelec.eu/areas-of-work/cen-cenelec-topics/artificial-intelligence/}{ISO/CEN-CENELEC JTC'21}}
    {\small \href{mailto:christos.hadjinikolis@gmail.com}{christos.hadjinikolis@gmail.com}}

\begin{aside}
    \section{prof\/iles}
        {\footnotesize \profilelink{profileweb}{https://christos-hadjinikolis.github.io/}{website}\\\profilelink{profilelinkedin}{https://uk.linkedin.com/in/christoshadjinikolis}{LinkedIn}\\\profilelink{profilegithub}{https://github.com/Christos-Hadjinikolis}{GitHub}\\\profilelink{profilescholar}{https://scholar.google.com/citations?hl=en&user=Rr3KTx8AAAAJ}{Google Scholar}}
    \section{leadership}
        {\footnotesize 6 direct reports\\career/performance\\hiring/coaching\\delivery ownership}
    \section{connectivity}
        {\footnotesize tool routing\\approval gates\\eval/replay\\auditable LLMs\\tool connectors}
    \section{platform}
        {\footnotesize distributed systems\\Kafka/Flink\\model/API serving\\observability\\incident response}
    \section{trust}
        {\footnotesize AI governance\\open standards\\auditability\\explainability\\enterprise trust}
    \section{open source}
        {\footnotesize \textbf{PyPI}: \href{https://pypi.org/project/dynamicio/}{dynamicio}\\\href{https://pypi.org/project/skeleton-replay/}{skeleton-replay}\\\href{https://plugins.jetbrains.com/plugin/32807-skeleton-replay}{JetBrains plugin}}
\end{aside}

\section{about me}
{\normalsize Engineering Manager with 16 years across backend platforms, production ML, data systems, and AI research. Manages \textbf{6} direct reports and leads cross-functional teams responsible for reliable, observable model-backed services. Recent work spans streaming platforms, model and API serving, eval/replay loops, fallback paths, incident response, and AI governance; Promet adds private hands-on GenAI work on tool routing, approval gates, traces, and auditable LLM workflows.}
\section{work-experience}
\begin{entrylist}
    \entry
    {12/2020--present}
    {Vortexa Ltd, London, UK}
    {Engineering Manager | ML Systems Lead}
    {
        {\normalsize
        \setlength{\parindent}{0pt}
        \setlength{\leftskip}{0pt}
        Own engineering strategy and delivery for a live post-ingestion ML/data estate turning \textbf{~6M} vessel-position records/hour into production intelligence for \textbf{13.5K} monitored vessels. Manage \textbf{6} direct reports and lead a \textbf{10-person} cross-functional team accountable for model/API serving, data quality, stakeholder alignment, reliability, and delivery maturity.

        }
        \begin{itemize}[leftmargin=1.1em]
            \item \normalsize\textbf{People \& Delivery:} Manage \textbf{6} direct reports across Product, SME analysis, Data Science, and Data Engineering; own 1:1s, performance/career development, hiring for \textbf{6} roles, coaching, delivery accountability, and team cadence.
            \item \normalsize\textbf{Platform \& Reliability:} Lead a distributed \textbf{Kafka/Flink} streaming estate; maintain rollback/fallback paths, shared on-call, incident response, observability/runbooks, and \textbf{MTTR} under \textbf{30} min.
            \item \normalsize\textbf{Eval/Replay \& Serving:} Led 0-to-1 research-to-production delivery using \textbf{PyTorch} for destination/arrival-time sequence and transformer models; established ML{f}low, model/data versioning, model/API serving, evaluation gates, replay, and monitoring.
            \item \normalsize\textbf{Trust, Interfaces \& Governance:} Drove data-lineage, schema-management, Avro/interface-contract, SLA, and domain-owner proposals so cross-team dependencies were explicit and governable.
            \item \normalsize\textbf{Product / Stakeholder Alignment:} Ran Product/SME/analyst/engineer workshops to turn model-quality disputes into shared definitions, measurable objectives, interface proposals, and OKR-linked workstreams.
            \item \normalsize\textbf{Engineering Ecosystem:} Standardised I/O seams, repo archetypes, AWS CodeArtifact/shared-library publishing, ADRs, dev containers, and local E2E tests (cut a \textbf{3h} DAG feedback loop to under \textbf{5} min).
        \end{itemize}
    }
    \entry
    {04/2016--12/2020}
    {DataReply, London, UK}
    {Senior Consultant | Data Scientist / ML Engineer}
    {
        {\normalsize
        \setlength{\parindent}{0pt}
        \setlength{\leftskip}{0pt}
        Joined Data Reply's London spin-off as its first consultant and helped grow it from \textbf{three practitioners and two managers to 30+ consultants}. Progressed to Senior Consultant; shaped enterprise client opportunities, timelines, expectations, consultant placement, and technical roadmaps across UBS, Vodafone, CNHi, and other clients.

        }
        \begin{itemize}[leftmargin=1.1em]
            \item \normalsize\textbf{Vodafone Group:} Led multi-workstream GCP data-science platform (based on Kubeflow) delivery; drove requirements, migration planning, CI/CD, automation, and ML lifecycle support.
            \item \normalsize\textbf{CNHi:} Lead Data Scientist/Scrum Master on predictive maintenance (Azure/Databricks); translated stakeholder needs into DS/DE backlog and guided a \textbf{PySpark} team.
            \item \normalsize\textbf{UBS:} Embedded consultant who grew from Data Scientist into ML Engineer through XP/pairing while delivering graph/process-mining analytics, observability, and real-time \textbf{Kafka}/Elasticsearch pipelines.
        \end{itemize}
    }
    \entry
    {2010--2016}
    {KCL, UCL, David Game College, UK}
    {Teaching and Technical Training}
    {
        {\small
        \setlength{\parindent}{0pt}
        \setlength{\leftskip}{0pt}
        Delivered courses and tutorials in algorithm design, data structures, databases/SQL, and Java/Python during and after doctoral research; developed presentation discipline, learner support, and mentoring habits.

        }
    }
\end{entrylist}
\newpage
\setlength{\parskip}{2pt}
\setlength{\parsep}{0pt}
\specialsection{education}
\begin{entrylist2}
    \entrytwo
        {\small 2010--2014}
        {\small Ph.D. in Computer Science}
        {\small King's College London, UK}
        {\small \emph{\textbf{Thesis}: Persuasion Dialogues \& Opponent Modelling.}

        {\footnotesize Developed logical inference methods for large \textbf{Knowledge Graphs}, combining symbolic reasoning, \textbf{graph analytics}, and Bayesian techniques (e.g., MCMC). Received \textbf{Best Poster Award} at IJCAI 2013 (413 submissions), \textbf{EPSRC PhD Scholarship}, \textbf{Outstanding TA Award}, and \textbf{Graduate Certif\/icate in Academic Practice}.
        \\ \textbf{\textcolor{blue}{\href{https://scholar.google.com/citations?hl=en&user=Rr3KTx8AAAAJ}{\underline{Access Publications on Google Scholar}}}}}
        }
    \entrytwo
        {\small  2004--2010}
        {\small Diploma (BEng) in Computer Engineering}
        {\small University of Thessaly, GR}
        {\footnotesize Five-year polytechnic degree in Computer Science and Engineering with a \textbf{strong focus on mathematics} (10 courses), covering theoretical aspects, networking, and telecommunications. \emph{\textbf{Majored in Artif\/icial Intelligence}}.}
    \entrytwo
        {\small  2002--2004}
        {\small Military Of\/f\/icer: Second Lieutenant}
        {\small National Guard, CY}
        {\small Completed national service as a commissioned Second Lieutenant.}
\end{entrylist2}
\specialsection{selected AI/connectivity tooling}
\begin{entrylist2}
    \entrytwo
        {\small Private}
        {\small \href{https://christos-hadjinikolis.github.io/2026/08/01/skeleton-replay-runtime-architecture-evidence.html}{Promet}}
        {\small GenAI harness}
        {\footnotesize Private applied GenAI project for memory-aware, voice-enabled, tool-using assistants; hands-on work on local runtimes, MCP-style connector patterns, tool routing, schema validation, approval gates, durable state/memory, traces, streaming UX, and deterministic replay/eval for auditable LLM interactions.}
    \entrytwo
        {\small 2026}
        {\small \href{https://christos-hadjinikolis.github.io/2026/05/18/mcp-is-plumbing-trust-is-the-product.html}{MCP: Plumbing, Trust}}
        {\small Public writing}
        {\footnotesize Practical article on MCP mental models, connector plumbing, staged tool exposure, host-side routing, approval, and evidence boundaries for local/private AI assistants.}
    \entrytwo
        {\small PyPI}
        {\small \href{https://pypi.org/project/skeleton-replay/}{skeleton-replay}}
        {\small Runtime evidence}
        {\footnotesize Python developer tool that turns script/pytest runs into traces, architecture snapshots, workflow evidence, and replayable reports for review, debugging, onboarding, and LLM-assisted code understanding.}
    \entrytwo
        {\small JetBrains}
        {\small \href{https://plugins.jetbrains.com/plugin/32807-skeleton-replay}{Skeleton Replay plugin}}
        {\small IDE workflow}
        {\footnotesize Free PyCharm/IntelliJ plugin bringing Skeleton reports, workflow evidence, and source navigation into the IDE; reflects a bias for developer experience, evidence, and shared understanding.}
    \entrytwo
        {\small PyPI}
        {\small \href{https://pypi.org/project/dynamicio/}{dynamicio}}
        {\small Data contracts}
        {\footnotesize Free public tooling for making I/O seams explicit, validating schema/data expectations, and switching local/dev/prod datasets so pipeline logic can be tested with faster feedback.}
\end{entrylist2}
\specialsection{standards/committees}
\begin{entrylist2}
    \entrytwo
        {\footnotesize since 10/2024}
        {\small University College London, Dept. of Information Studies}
        {\small Associate Researcher}
        {\small Research/teach practical AI standardisation: AI Act, auditability, model/data versioning, explainability, and safe adoption.}
    \entrytwo
        {\footnotesize since 01/2021}
        {\small ISO - International Organisation for Standardisation}
        {\small Committee Expert Member for JTC'21}
        {\footnotesize Member of JTC'21 Working Group 3, developing AI standards aligned with EU policy and international norms.}
\end{entrylist2}
\specialsection{talks/interviews}
\begin{entrylist2}
    \entrytwo
        {\small 2023}
        {\small Bill Raymonds, Agile in Action}
        {\small Podcast Interview}
        {\small \textbf{Agile Data Science}: Sailing Through Data Science: The Agile Journey at Vortexa, \underline{\textbf{Details:} \href{https://christos-hadjinikolis.github.io//2023/10/31/Agile-In-Action.html}{YouTube}}}
    \entrytwo
        {\small 2022}
        {\small Open Data Science Conference}
        {\small Industry Conference Talk}
        {\small \textbf{dynamicio}, a published \href{https://pypi.org/project/dynamicio/}{\textbf{PyPI library}} for abstracting I/O in ML/data workflows, \underline{\textbf{Details:} \href{https://odsc.com/speakers/dynamicio-a-pandas-i-o-wrapper-why-you-should-start-your-ml-ops-journey-with-wrapping-your-i-o/}{here}}}
    \entrytwo
        {\small 2020}
        {\small iunera: Big Data Innovators}
        {\small Interview Blog Post}
        {\small The Agile approach in data science explained by an ML expert, \underline{\textbf{Details:} \href{https://www.iunera.com/kraken/interviews/the-agile-approach-in-data-science-explained-by-an-ml-expert/}{here}}}
    \entrytwo
        {\small 2020}
        {\small Big Data Warsaw}
        {\small Industry Conference Talk}
        {\small Monitoring \& Analysing Communication and Trade Events as Graphs, \underline{\textbf{Details:} \href{https://bigdatatechwarsaw.eu/agenda/}{here}}}
    \entrytwo
        {\small 2018}
        {\small Connected Data London}
        {\small Panelist @ Industry Conference}
        {\small AI tribes and interconnections: what's graph got to do with it?, \underline{\textbf{Details:} \href{http://connected-data.london/2018/09/23/ai-tribes-and-interconnections-whats-graph-got-to-do-with-it/}{here}}}
    \entrytwo
        {\small 2018}
        {\small Minds Mastering Machines}
        {\small Industry Conference Talk}
        {\small A talk on doing Data Science the Agile Way, \underline{\textbf{Details:} \href{https://2018.mcubed.london/sessions/data-science-agile-way/}{here}}}
\end{entrylist2}
\specialsection{credentials \& awards}
\begin{entrylist2}
    \entrytwo
        {\small 2017--2020}
        {\small Selected ML/cloud credentials}
        {\small AWS, Google Cloud, O'Reilly}
        {\footnotesize AWS ML Specialty; Google Professional Data Engineer; Apache Spark Developer.}
    \entrytwo
        {\small 2016--2019}
        {\small Data/graph engineering credentials}
        {\small Coursera, Elastic, Neo4j}
        {\footnotesize Process Mining; Graph Analytics; Advanced Elasticsearch; Neo4j Certif\/ied Professional.}
\end{entrylist2}
\end{document}
